% Supplementary Materials: Dictionarium Sinicum
% Compile with LuaLaTeX
\documentclass[11pt,a4paper]{article}

\usepackage[review]{acl}
\usepackage{microtype}
\usepackage{booktabs}
\usepackage{float}  % for [H] placement

% Font setup for LuaLaTeX
\usepackage[english,bidi=default]{babel}
\babelfont{rm}{TeXGyreTermesX}

% CJK
\usepackage{luatexja-fontspec}
\setmainjfont{Noto Serif CJK SC}

% Prevent luatexja from inserting inter-character spacing around Cyrillic
\ltjsetparameter{jacharrange={-2}}

% Korean
\babelprovide[import]{korean}
\babelfont[*hangul]{rm}{Noto Serif CJK KR}

% Thai
\babelprovide[import]{thai}
\babelfont[*thai]{rm}{Noto Sans Thai}

% Devanagari (Hindi)
\babelprovide[import]{hindi}
\babelfont[*devanagari]{rm}{Noto Sans Devanagari}

% Russian / Cyrillic — direct font, babel doesn't switch in tables
\newfontfamily\cyrfont{Noto Serif}

% Arabic — babel with proper bidi support
\babelprovide[import]{arabic}
\babelfont[*arabic]{rm}{Noto Sans Arabic}

\title{Supplementary Materials:\\Dictionarium Sinicum}
\author{Anonymous}

% Helper commands for script switching inside tables
\newcommand{\ru}[1]{{\cyrfont #1}}
\newcommand{\ar}[1]{\foreignlanguage{arabic}{#1}}
\newcommand{\hi}[1]{\foreignlanguage{hindi}{#1}}
\newcommand{\thi}[1]{\foreignlanguage{thai}{#1}}

% Allow multiple table* floats per page
\setcounter{dbltopnumber}{3}
\renewcommand{\dbltopfraction}{0.9}

\begin{document}
\maketitle

\section{Sample Dictionary Entries}

We present three sample entries spanning common vocabulary, concrete nouns, and culturally significant loanword sources.
Each entry shows definitions from community dictionaries and LLM-generated glosses across 18 languages.

\subsection{你好 [nǐ hǎo] --- Common Greeting}

\begin{table}[H]
\centering
\footnotesize
\begin{tabular}{llp{3.6cm}}
\toprule
\textbf{Lang} & \textbf{Source} & \textbf{Definition} \\
\midrule
en & cedict & hello; hi \\
de & handedict & Hallo! / Guten Tag! \\
de & minimax & Hallo!; Guten Tag! \\
fr & cfdict & bonjour / salut \\
fr & minimax & bonjour; salut \\
es & chedicc & hola; ¿qué tal? \\
id & cidict & halo; hai \\
ja & jmdict & ニーハオ \\
ja & minimax & こんにちは \\
ko & minimax & 안녕; 안녕하세요 \\
ru & minimax & \ru{привет} \\
vi & minimax & xin chào; chào \\
tl & minimax & hello; kamusta \\
sv & minimax & hej; hallå \\
nl & minimax & hallo; dag \\
pt & minimax & olá; oi \\
it & minimax & ciao; buongiorno \\
ar & minimax & \ar{مرحبا} / \ar{أهلا} \\
th & minimax & \thi{สวัสดี} \\
hi & minimax & \hi{नमस्ते}; \hi{हेलो} \\
fa & minimax & \ar{سلام} \\
\bottomrule
\end{tabular}
\caption{Dictionary entry for 你好 \textit{nǐ hǎo} (hello). All 18 target languages: community sources and MiniMax glosses spanning Latin, CJK, Cyrillic, Arabic, Thai, and Devanagari scripts.}
\label{tab:nihao}
\end{table}

\begin{table}[H]
\centering
\footnotesize
\begin{tabular}{llll}
\toprule
\textbf{Headword} & \textbf{Meaning} & \textbf{Cantonese} & \textbf{Hokkien} \\
\midrule
你好 & hello    & \textit{nei5 hou2}   & \textit{lí hó} \\
银行 & bank     & \textit{ngan4 hong4}  & \textit{gîn-hâng} \\
豆腐 & tofu     & \textit{dau6 fu6}     & \textit{tāu-hū} \\
\bottomrule
\end{tabular}
\caption{Dialect forms for the three sample entries. Cantonese in Jyutping romanization, Hokkien in POJ.}
\label{tab:dialect-samples}
\end{table}

\subsection{银行 [yín háng] --- Bank}

\begin{table}[H]
\centering
\footnotesize
\begin{tabular}{llp{3.6cm}}
\toprule
\textbf{Lang} & \textbf{Source} & \textbf{Definition} \\
\midrule
en & cedict & bank \\
en & wiktextract & bank (financial institution); jewellery dealer \\
de & handedict & Bank (S, Wirtsch) \\
fr & cfdict & banque \\
id & cidict & bank \\
ja & jmdict & ぎんこう \\
ko & minimax & 은행; 금융 기관 \\
es & minimax & banco \\
ru & minimax & \ru{банк} \\
vi & minimax & ngân hàng \\
tl & minimax & bangko \\
th & minimax & \thi{ธนาคาร} \\
ar & minimax & \ar{مصرف} / \ar{بنك} \\
hi & minimax & \hi{बैंक} \\
\bottomrule
\end{tabular}
\caption{Dictionary entry for 银行 (bank). Community sources + selected MiniMax languages.}
\end{table}

\subsection{豆腐 [dòu fu] --- Tofu (Loanword Source)}

This entry illustrates both the multilingual definitions and the dialect forms that enable loanword tracing across Southeast Asia.

\begin{table}[H]
\centering
\footnotesize
\begin{tabular}{llp{3.6cm}}
\toprule
\textbf{Lang} & \textbf{Source} & \textbf{Definition} \\
\midrule
en & cedict & tofu; bean curd \\
de & handedict & Tofu, Doufu, Bohnenkäse (S, Ess) \\
fr & cfdict & tofu \\
id & cidict & tahu; kembang tahu \\
ja & jmdict & とうふ \\
ko & minimax & 두부 \\
tl & minimax & tokwa \\
vi & minimax & đậu phụ \\
th & minimax & \thi{เต้าหู้} \\
\bottomrule
\end{tabular}
\caption{Dictionary entry for 豆腐 (tofu).}
\end{table}

\noindent The Hokkien form \textit{tāu-hū} (Table~\ref{tab:dialect-samples}) is the etymological source for Indonesian \textit{tahu}, Tagalog \textit{tokwa} (via Hokkien \textit{tāu-koa\textsuperscript{n}} ``dried tofu''), and Thai \thi{เต้าหู้}---all visible in the multilingual definitions above.

%% table* floats declared here (during page 1) so they float to page 2 top
%% and get numbered 5--7 (after the page 1 [H] tables 1--4)

%% MT baseline — full-width, 3 metric groups side-by-side
\begin{table*}[t]
\centering
\footnotesize
\begin{tabular}{lrccc@{\hskip 1.5em}rccc@{\hskip 1.5em}rccc}
\toprule
& \multicolumn{4}{c}{\textit{Sense coverage (\%)}} & \multicolumn{4}{c}{\textit{False sense rate (\%)}} & \multicolumn{4}{c}{\textit{Format compliance (\%)}} \\
\cmidrule(lr){2-5} \cmidrule(lr){6-9} \cmidrule(l){10-13}
\textbf{System} & \textbf{n} & \textbf{de} & \textbf{fr} & \textbf{id} & \textbf{n} & \textbf{de} & \textbf{fr} & \textbf{id} & \textbf{n} & \textbf{de} & \textbf{fr} & \textbf{id} \\
\midrule
Our pipeline      & 453 & 79.2 & 84.0 & 93.7 & 453 & 19.8 & 14.2 &  6.3 & 453 & ${\sim}$100 & ${\sim}$100 & ${\sim}$100 \\
MiniMax (generic)  & 177 & 58.3 & 67.8 & 83.6 & 177 & 45.4 & 29.1 & 17.6 & 177 & 97.9 & 86.4 & 89.5 \\
Claude Sonnet 4    & 275 & 61.1 & 64.6 & 87.9 & 275 & 47.2 & 37.0 & 15.2 & 275 & 79.5 & 71.4 & 74.1 \\
GPT-4o-mini        & 275 & 52.2 & 61.9 & 83.8 & 275 & 52.9 & 41.3 & 17.4 & 275 & 78.1 & 73.0 & 72.7 \\
Google Translate   & 275 & 52.4 & 60.1 & 85.5 & 275 & 51.6 & 36.5 & 13.5 & 275 & 76.7 & 74.6 & 68.3 \\
DeepL              & 275 & 49.3 & 52.4 & 83.9 & 275 & 54.4 & 47.2 & 15.7 & 275 & 75.3 & 68.3 & 68.3 \\
Azure Translator   & 275 & 48.7 & 60.3 & 86.4 & 275 & 53.0 & 39.1 & 12.9 & 275 & 72.6 & 68.3 & 68.3 \\
\bottomrule
\end{tabular}
\caption{Per-language MT baseline breakdown across three metrics. Our pipeline's advantage is most pronounced on German, where baselines score 20--30 points lower on sense coverage. Format compliance measures adherence to CEDICT-style gloss lists.}
\label{tab:mt-detail}
\end{table*}

%% Back-translation + Judge — side by side
\begin{table*}[t]
\centering
\footnotesize
\begin{tabular}{lrcc@{\hskip 3em}lrccc}
\toprule
\multicolumn{4}{c}{\textit{Back-translation evaluation (\%)}} & \multicolumn{5}{c}{\textit{LLM-as-judge scores (1--5)}} \\
\cmidrule(r){1-4} \cmidrule(l){5-9}
\textbf{Language} & \textbf{n} & \textbf{Sense} & \textbf{False} &
\textbf{Language} & \textbf{n} & \textbf{Acc.} & \textbf{Nat.} & \textbf{Comp.} \\
\midrule
Spanish    & 126 & 79.0 & 20.2 & Russian    &  6 & 4.75 & 4.50 & 3.75 \\
Portuguese & 123 & 75.2 & 22.3 & Arabic     &  3 & 4.67 & 4.50 & 4.00 \\
Italian    & 122 & 74.7 & 25.4 & Spanish    &  8 & 4.56 & 4.44 & 4.00 \\
Swedish    & 118 & 74.2 & 24.2 & Dutch      & 12 & 4.54 & 4.58 & 3.58 \\
Russian    & 123 & 71.3 & 29.4 & Portuguese & 10 & 4.50 & 4.50 & 3.35 \\
Dutch      & 125 & 69.7 & 28.2 & Vietnamese & 10 & 4.50 & 4.30 & 3.60 \\
Persian    & 128 & 67.6 & 31.6 & Swedish    &  6 & 4.42 & 4.58 & 3.50 \\
Thai       & 113 & 66.8 & 31.4 & Persian    &  8 & 4.25 & 4.56 & 3.00 \\
Vietnamese & 107 & 65.4 & 28.2 & Korean     &  6 & 4.25 & 4.33 & 2.92 \\
Arabic     & 129 & 65.3 & 34.7 & Italian    &  6 & 4.00 & 4.08 & 2.92 \\
Hindi      & 128 & 63.4 & 32.3 & Thai       &  7 & 4.07 & 4.36 & 3.14 \\
Tagalog    & 110 & 61.7 & 34.1 & Hindi      & 13 & 4.04 & 3.77 & 3.04 \\
Korean     & 119 & 60.0 & 38.2 & Tagalog    &  5 & 3.90 & 3.60 & 3.60 \\
\midrule
\textbf{Overall} & \textbf{1{,}571} & \textbf{68.9} & \textbf{29.2} &
\textbf{Overall} & \textbf{100} & \textbf{4.34} & \textbf{4.31} & \textbf{3.38} \\
\bottomrule
\end{tabular}
\caption{Evaluation of 13 non-community languages. Left: back-translation sense coverage. Right: LLM-as-judge quality scores (accuracy, naturalness, completeness). Judges: Claude Sonnet~4, Kimi~K2, GPT-OSS-120B, DeepSeek~V3.1. Mean pairwise inter-judge disagreement: 0.64~points.}
\label{tab:backtrans-judge}
\end{table*}

%% Eval scripts — full-width
\begin{table*}[t]
\centering
\footnotesize
\begin{tabular}{lp{11cm}}
\toprule
\textbf{Script} & \textbf{Purpose} \\
\midrule
\texttt{eval\_community\_comparison.py} & Core evaluation against 4 community dictionaries (de, fr, id). Shared stratified sampling and metrics. \\
\texttt{eval\_mt\_baseline.py} & MT baselines: MiniMax generic prompt, Google Translate, DeepL, Azure, Claude Sonnet~4, GPT-4o-mini. \\
\texttt{eval\_multi\_model.py} & Free LLM comparison via Groq, Cerebras, and SambaNova inference providers. \\
\texttt{eval\_backtranslation.py} & Back-translation round-trip + LLM-as-judge quality scoring for 13 non-community languages. \\
\bottomrule
\end{tabular}
\caption{Evaluation script manifest. All require Python~3.11+ and a populated \texttt{dictmaster.db}. Results in \texttt{eval-results/} (JSON).}
\label{tab:scripts}
\end{table*}

%% ============================================================
\section{Evaluation Details}

Tables~\ref{tab:mt-detail}--\ref{tab:scripts} present the full evaluation results and script manifest.
All raw results are provided as JSON in \texttt{eval-results/}.

\smallskip
\noindent\textbf{Sampling.}
Headwords are stratified by frequency (high: $\geq$100 articles, mid: 10--99, low: 1--9, rare: 0), 50 per band, \texttt{seed=42}, yielding ${\sim}$275 evaluation pairs across de/fr/id.

\smallskip
\noindent\textbf{Reproducibility.}
MiniMax M2.5 (\texttt{abab7-chat-preview}) via \texttt{api.minimax.chat}; \$0.30/M input, \$1.20/M output.
Database: 428K headwords $\times$ 18 languages in SQLite.
MT baselines: Google Translate v2, DeepL Free, Azure Translator v3, Claude Sonnet~4, GPT-4o-mini; free LLMs via Groq, Cerebras, SambaNova.

\end{document}
